\chapter{SYSTEM ANALYSIS AND REQUIREMENTS}
\thispagestyle{empty}

This chapter records the functional and non-functional requirements that the
implemented system satisfies, describes the user roles, the system
constraints, and the software and hardware environment.

\clearpage

This chapter records the requirements that the implemented system satisfies,
derived from the approved proposal and verified against the repository.

\section{Functional Requirements}

\begin{enumerate}
  \item \textbf{Identity and tenancy.} Users register and log in with
        cookie-based JWT sessions that rotate refresh tokens and include CSRF
        double-submit protection (\textsc{fr-1}).
  \item \textbf{Institute membership.} A user belongs to institutes through
        memberships, each carrying roles, and every request requires an
        institute selector header that is validated against an active
        membership (\textsc{fr-2}).
  \item \textbf{Role based access.} Actions are restricted to the roles
        \textsc{institute\_admin}, \textsc{teacher}, and \textsc{student}
        (\textsc{fr-3}).
  \item \textbf{Academic structure.} Subjects, chapters, and topics can be
        created, edited, soft-deleted, restored, and reassigned, with unique
        slugs per parent (\textsc{fr-4}).
  \item \textbf{Syllabus automation.} A syllabus document can be proposed as
        a chapter--topic hierarchy by an AI job; the proposal remains pending
        review until a teacher or administrator confirms it, after which the
        hierarchy becomes authoritative (\textsc{fr-5}).
  \item \textbf{Materials and OCR.} PDF, image, and text materials are
        uploaded (with chunked upload for large files), processed through the
        OCR service, retried on failure, and cancellable; per-page inspection
        and correction are supported (\textsc{fr-6}).
  \item \textbf{Material Intelligence.} An enhanced, versioned view of a
        material is produced by cleaning, structuring, and segmenting the
        extracted text and mapping each segment to syllabus entities with a
        confidence value (\textsc{fr-7}).
  \item \textbf{Study content.} AI jobs generate notes, summaries,
        flashcards, and concepts with type-specific validated payloads and
        append-only version history (\textsc{fr-8}).
  \item \textbf{Question domain.} The question bank supports multiple
        question types, difficulty levels, teacher review, AI batch
        generation, and deterministic extraction of question candidates with
        explicit ambiguity reports (\textsc{fr-9}).
  \item \textbf{Paper patterns.} Nested pattern structures with per
        question-type attempt rules are validated, approved, locked, and
        extractable from text, PDF, and image sources (\textsc{fr-10}).
  \item \textbf{Question papers and assessments.} Papers are authored against
        a pattern with coverage gating, AI shortfall-generation, scoped
        question selection, publication, and export to PDF, DOCX, and Excel
        formats (\textsc{fr-11}).
  \item \textbf{Examinations.} Assessments are scheduled; attempts snapshot
        the question set immutably, enforce deadlines server-side, validate
        per-type answers, evaluate objective question types deterministically,
        and produce result analytics (\textsc{fr-12}).
  \item \textbf{Practice.} Flashcard and question practice sessions provide
        per-item instant feedback without scoring (\textsc{fr-13}).
  \item \textbf{Export.} Content, questions, assessments with results, paper
        patterns, and question papers are exported to DOCX, PDF, and spreadsheets
        (\textsc{fr-14}).
  \item \textbf{Job lifecycle.} Long-running operations run as asynchronous
        jobs with states \textsc{queued}, \textsc{processing},
        \textsc{completed}, \textsc{failed}, and \textsc{cancelling}/
        \textsc{cancelled}, with retry and cancel endpoints (\textsc{fr-15}).
\end{enumerate}

\section{Non-Functional Requirements}

\begin{enumerate}
  \item \textbf{Tenant isolation.} Every query that touches user data is
        scoped to an institute; isolation is enforced in application code with
        tenant-scoped filters and foreign keys (\textsc{nfr-1}).
  \item \textbf{Security.} Session cookies are httpOnly; refresh tokens are
        stored hashed and rotated with race handling; CSRF protection covers
        all state-changing requests; internal services authenticate with an
        internal API key; answer keys and explanations never leave the API
        (\textsc{nfr-2}).
  \item \textbf{Responsiveness.} The API never blocks on long-running
        operations; heavy computation is delegated to async workers
        (\textsc{nfr-3}).
  \item \textbf{Reliability of AI output.} AI output is validated against
        canonical schemas on both worker and API boundaries with automatic
        retry on validation failure (\textsc{nfr-4}).
  \item \textbf{Maintainability.} The codebase is a typed monorepo with linting,
        type checking across workspaces, and per-layer unit tests
        (\textsc{nfr-5}).
  \item \textbf{Reproducibility and documentation.} Configuration is
        environment-driven with template files; task tracking, project status,
        architecture documents, and API reference are maintained in the
        repository (\textsc{nfr-6}).
\end{enumerate}

\section{User Roles}

\begin{description}[leftmargin=1em]
  \item[\textsc{Institute Administrator}]\hfill\\
        Manages institute members and their roles, subjects, and platform
        settings; may confirm syllabus proposals and approve patterns.
  \item[\textsc{Teacher}]\hfill\\
        Creates subjects, chapters, and topics; uploads and processes
        materials; reviews and confirms AI proposals; manages the question
        bank, paper patterns, question papers, and assessments; reviews
        results and analytics.
  \item[\textsc{Student}]\hfill\\
        Views available assessments, attempts them within the schedule with a
        timer and deadline enforcement, submits, and views evaluated results;
        practices with flashcards and questions.
\end{description}

\section{System Constraints}

\begin{enumerate}
  \item The main application is a modular monolith; no premature
        microservices. The OCR service is the only separately deployable
        component.
  \item PostgreSQL is the only primary database; JSONB is used for flexible
        structured content.
  \item RabbitMQ is used for asynchronous task distribution; long-running
        work never runs inside the API process.
  \item All AI access flows through an internal OpenAI-compatible gateway
        (OmniRoute); no local LLM and no direct cloud provider SDK calls are
        used by the API or workers.
  \item The Cloudflare Tunnel is the only public network ingress; internal
        services are not exposed on host ports in the production compose file.
\end{enumerate}

\section{Software and Hardware Environment}

\subsection{Software environment}

\begin{itemize}
  \item Operating environments: Linux containers under Docker Compose;
        development on Linux machines.
  \item Runtimes: Node.js (TypeScript, strict, ES2022) for the API and web
        applications; Python 3.12 for the OCR service and workers.
  \item Services: PostgreSQL 17, RabbitMQ 3, an internal OpenAI-compatible AI
        gateway, nginx as reverse proxy, and cloudflared for the public
        tunnel.
  \item Build tooling: pnpm workspaces, Turborepo, Docker with cached
        dependency-first builds.
\end{itemize}

\subsection{Hardware environment}

Requirements are covered by the distributed design rather than a single
large host: the OCR service benefits from a GPU-compatible host but falls
back to CPU-based inference, and the standalone worker model allows OCR to be
offloaded to additional devices connected through the API. No specialized
hardware is required by the platform itself.

\section{Assumptions}

\begin{enumerate}
  \item Source documents are provided as PDF, common image formats, or plain
        text.
  \item Institutions provide syllabus and material documents that AI and
        deterministic extractors can process; ambiguous sources produce
        explicit ambiguity reports that a human resolves.
  \item Teachers review and confirm AI-proposed structure before it becomes
        authoritative.
  \item An OpenAI-compatible endpoint is available through the internal
        gateway for AI generation.
\end{enumerate}